\honest{Hand vs. Fingers}{Eliezer Yudkowsky}{2008}

Back to reductionism. If things had to be fundamental to be fun, we could take joy in nothing bigger than a quark, so I reject that view. The thesis, again: we use models with many levels because they are cheaper to compute, but reality has only one level.

When you pick up a cup, is it your hand that picks it up? A joke: scientists have measured the forces of your fingers, thumb and palm and found ``no force left over,'' so the force of your hand must be zero. \nb{The joke has the form of a real argument in philosophy of mind, the exclusion problem: if a physical property does all the causal work, what is left for a mental one to do? The post's answer, that the hand is its parts, is one of the standard replies.}

My theme is that once you can see how a higher level reduces to a lower one, the two stop seeming like separate things, and it becomes silly to argue about whether the hand or the fingers pick up the cup. ``See'' means concrete visualization: imagining a hand makes you imagine fingers, thumb and palm, and imagining those makes you find a hand in the picture, so the two levels of your map are bound tightly together in your mind. In reality the levels are bound together by ``the tightest possible binding: physical identity.'' ``Hand'' and ``fingers and thumb and palm'' name one thing from two points of view. \nb{Philosophers call this view composition as identity, and the \textsc{Stanford Encyclopedia of Philosophy} calls it ``by no means an uncontroversial view.'' The argument needs only that the hand is made of its parts.}

Suppose instead you had a hand scanner that showed the hand as a dot, like an old radar screen, and other scanners for the fingers, thumb and palm. You would see a cluster of dots, but you could imagine the hand-dot moving off on its own, though with the sensors as they are that is physically impossible. Even if told, or if you guessed from the dots, you would only know the reduction, not see it. If people merely told you ``There's a hand over there, and some fingers over there,'' you would be like an old-fashioned AI with LISP tokens: it could assert Inside(Room, Hand) and not-Inside(Room, Fingers) with no contradiction, lacking the rule that whatever contains the hand contains the fingers. None of this lets a hand crawl ghostlike across the room. It is a fact about the map. What seems possible to a particular person is not what is logically or physically possible. It may seem possible to you that 235757 is prime, but it is not; a logically omniscient mind would see the factors. That is why we have impossible possible worlds, so that we can put probabilities on propositions that may in fact be logically impossible. \nb{It is $431 \times 547$.}

Then I imagine philosophers who attack ``eliminative fingerists'' for denying the direct fact that we feel our hand hold the cup, since if hands did not exist the cup would fall, and philosophers who propose ``appendigital bridging laws'' by which fingers call a hand into being, laws that could have been different. All of them commit the Mind Projection Fallacy: they take their intuitions for information about reality. \nb{Both parodies move views about the mind onto hands. Eliminative materialists deny that beliefs exist because they expect beliefs not to reduce to brain states. David Chalmers, a dualist, posits laws linking experience to the physical, and writes that ``it seems that we cannot coherently conceive'' of a world physically like ours without water. Whether the mind is like the hand is the question in dispute, and the post does not argue it.} Your inability to imagine something is a fact about your brain. ``Another brain might work differently.''
